\documentclass[11pt,a4paper]{report}
\usepackage[utf8]{inputenc}
\usepackage[T1]{fontenc}
\usepackage[french]{babel}
\usepackage{graphicx}
\usepackage{hyperref}
\usepackage{listings}
\usepackage{xcolor}
\usepackage{geometry}

\geometry{margin=2.5cm}

% Configuration pour les listings de code
\lstset{
  basicstyle=\ttfamily\small,
  breaklines=true,
  frame=single,
  numbers=left,
  numberstyle=\tiny,
  keywordstyle=\color{blue},
  commentstyle=\color{gray},
  stringstyle=\color{red}
}

\title{Rapport de Projet\\Programmation Internet\\Système de Partage de Fichiers P2P}
\author{Guillaume Xue : 22101031 \\ David Yu : 22110478}
\date{\today}

\begin{document}

\maketitle

\chapter{Introduction}

Ce projet implémente un système de partage de fichiers P2P utilisant UDP, un serveur central pour la découverte des pairs et un arbre de Merkle pour garantir l'intégrité des données.

\section{Objectifs}

\begin{itemize}
    \item S'enregistrer auprès d'un serveur central
    \item Partager et télécharger des fichiers via P2P
    \item Traverser les NAT avec relais
    \item Gérer efficacement les transferts avec fenêtre glissante
\end{itemize}

\section{Architecture}

L'implémentation en Go s'articule autour de : Transport UDP, Merkle (hashing fichiers), Crypto (ECDSA), Peer Manager, Protocol.

\section{Enregistrement au serveur}

\subsection{Clé publique et adresse IP}
La clé publique ECDSA est envoyée en PUT au serveur central. Le client établit un socket UDP dual-stack (IPv6 avec fallback IPv4) et effectue un handshake Hello/HelloReply avec le serveur.

\subsection{Keep-Alive}
Un mécanisme de keep-alive envoie périodiquement des pings (tous les 3min) au serveur central et à tous les pairs connectés pour maintenir les mappings NAT et détecter les pairs inactifs.

\section{Partage et téléchargement}

\subsection{Arbre de Merkle}
Chaque fichier/répertoire est représenté par un arbre de Merkle (SHA-256). Types : Chunk (données brutes, 1024 max), Directory (16 entrées), Big (32 chunks), BigDirectory (32 répertoires).

\subsection{Téléchargement avec fenêtre glissante}
Le downloader utilise 3 workers : SenderLoop (envoie requêtes), ResponseLoop (traite réponses), MonitorLoop (gère timeouts). La fenêtre s'ajuste dynamiquement en fonction des performances réseau.

\section{Fonctionnalités avancées}

\subsection{NAT Traversal}
Les clients derrière NAT peuvent communiquer via un relais (pair avec extension \texttt{NatTraversalRelay}). Processus : 
\begin{enumerate}
    \item Source demande au relais
    \item Relais notifie la cible
    \item Les deux exécutent du hole punching simultané avec backoff exponentiel (0s, 1s, 2s)
    \item NAT crée des mappings permettant la connexion directe
\end{enumerate}

\subsection{Autres fonctionnalités}
Interface CLI interactive, support natif IPv4/IPv6 (dual-stack avec fallback), PeerManager gère les adresses multiples de chaque pair, keep-alive multi-protocoles.

\section*{Crédits}

Ce projet a été réalisé avec l'aide des ressources suivantes :

\begin{itemize}
    \item \textbf{Documentation officielle} : Go standard library documentation
    \begin{itemize}
        \item \texttt{sync} : https://golang.org/pkg/sync/
        \item \texttt{time} : https://golang.org/pkg/time/
    \end{itemize}
    \item \textbf{Cours Université Paris Cité} : Programmation Réseau(mutex), Programmation Orientée Objet(structs, interfaces)
    \item \textbf{Ressources externes} : 
    \begin{itemize}
        \item \texttt{Merkle tree} : https://en.wikipedia.org/wiki/Merkle\_tree
    \end{itemize}
\end{itemize}

\end{document}
